\documentclass[11pt]{article}
\usepackage[margin=1in]{geometry}
\usepackage{amsmath,amssymb,amsthm}
\usepackage{microtype}
\usepackage[hidelinks]{hyperref}
\newtheorem{theorem}{Theorem}
\newtheorem{lemma}{Lemma}
\newcommand{\R}{\mathbb R}
\newcommand{\C}{\mathbb C}
\newcommand{\ind}{\mathbf 1}
\newcommand{\Mod}{\operatorname{Mod}_2}
\newcommand{\Adm}{\operatorname{Adm}}
\newcommand{\length}{\mathcal L^1}
\title{Infinite modulus between a compact real set and its real complement}
\author{Alec Kriebel\\
  \small ORCID: \href{https://orcid.org/0009-0001-9320-500X}{0009-0001-9320-500X}}
\date{17 October 2026}
\begin{document}
\maketitle

\begin{abstract}
Let $E\subset[0,1]$ be compact and let $F=\R\setminus E$ be its complement
within the real axis. We prove that positive planar $2$-modulus of the family
joining $E$ to $\{|z|=2\}$ implies infinite modulus of the family joining $E$
to $F$. Admissibility is imposed on every curve for pointwise extended Borel
densities. A weighted-distance potential and a critical trace estimate first
force any hypothetical finite-energy separator to have $\length(E)=0$; a
separate small-circle argument then treats every such null set, including
null sets of positive conformal capacity. The result is stated under explicit
conventions for the question recorded as Hayman--Lingham Problem~7.58.
\end{abstract}

\section{Statement and conventions}

We identify $\R$ with the horizontal axis in $\C=\R^2$, write $dA$ for
planar Lebesgue measure, and put $S_2=\{z:|z|=2\}$. A connector from $A$ to
$B$ is a continuous curve $\gamma:[0,1]\to\R^2$ with
$\gamma(0)\in A$, $\gamma(1)\in B$, locally rectifiable on $(0,1)$.
Equivalently, one may use the open parameter interval and require the stated
endpoint limits. For a nonnegative Borel function $\rho$, its line integral
is the increasing limit of the integrals on compact interior subcurves.
Arclength is counted with multiplicity; constant portions have zero length.
In particular, an infinite value at an endpoint alone contributes no mass.
The integral may be infinite, and the whole curve need not have finite
Euclidean length.

Let $\Gamma(A,B)$ consist of \emph{all} such connectors in the whole plane.
There is no condition that their interiors avoid either plate. Curves may
cross or revisit $A$, $B$, or the real axis before their terminal endpoints.
Define
\[
 \begin{split}
 \Mod\Gamma&=\inf_{\rho\in\Adm(\Gamma)}\int_{\R^2}\rho^2\,dA,\\
 \Adm(\Gamma)&=\bigl\{\rho:\R^2\to[0,\infty]\ \text{Borel}:\\
 &\hspace{3em}\int_\gamma\rho\,ds\geq1\quad\text{for every }\gamma\in\Gamma\bigr\}.
 \end{split}
\]
The infimum of an empty set is $\infty$. Densities are pointwise functions,
not equivalence classes modulo area-zero sets, and may be infinite on such
sets. We also prove the theorem when $\Gamma(A,B)$ is defined using only
globally rectifiable connectors, provided that convention is used
consistently in both modulus families.

\begin{theorem}\label{thm:main}
For every compact $E\subset[0,1]$, with $F=\R\setminus E$, one has
\[
 \Mod\Gamma(E,S_2)>0\quad\Longrightarrow\quad
 \Mod\Gamma(E,F)=\infty.
\]
\end{theorem}

The complement here is the real complement, not the planar complement.
Hayman and Lingham record the question, attributed to A.~A.~Gonchar and
communicated by M.~Vuorinen, in Problem~7.58 of the 2018 version of their
collection~\cite{HL}. The theorem specifies the curve and density conventions
used here; no conclusion about other interpretations of unspecified source
terminology is intended.

We prove the contrapositive. The intermediate conclusion $\length(E)=0$
is not sufficient: a compact null subset of the real axis can have positive
conformal capacity. Section~\ref{sec:circles} supplies the additional argument
at every endpoint of $E$.

\section{A pointwise separating potential}

The following majorization is pointwise, including at infinite values.

\begin{lemma}\label{lem:majorant}
If $Q$ is a compact rectangle and $r:Q\to[0,\infty]$ is Borel with
$\int_Qr^2\,dA<\infty$, there is a relatively lower semicontinuous
$g:Q\to[1,\infty]$ such that $g\geq r$ everywhere and
$\int_Qg^2\,dA<\infty$.
\end{lemma}
\begin{proof}
Set $h=r^2$ and $A_k=\{h>2^k\}$, $k\geq0$. By outer regularity choose
relatively open $O_k\supset A_k$ with
$|O_k|\leq|A_k|+2^{-2k-4}$, where $|\cdot|$ denotes planar area in this
proof. Put
\[
 H=1+\sum_{k\geq0}2^{k+1}\ind_{O_k},\qquad g=\sqrt H.
\]
Open-set indicators, finite sums, and increasing limits are lower
semicontinuous. Thus $g$ is lower semicontinuous. Dyadic geometric sums give
$H\geq h$ at every point: the constant term covers $h\leq1$, and if
$h=\infty$, the point belongs to every $O_k$. Moreover,
$\sum_{k\geq0}2^{k+1}\ind_{\{h>2^k\}}\leq4h$ for finite $h\geq0$.
Tonelli therefore gives
\[
 \int_Qg^2\,dA\leq |Q|+4\int_Qr^2\,dA+
       \sum_{k\geq0}2^{-k-3}
 =|Q|+4\int_Qr^2\,dA+\tfrac14<\infty.
\]
In particular, any infinite values retained by $g$ form an area-zero set.
\end{proof}

Fix $Q=[-2,3]\times[-2,2]$. For a nonempty compact $E\subset[0,1]$ and
a lower semicontinuous $g\geq1$ on $Q$ with finite squared integral, define
\begin{equation}\label{eq:distance}
 d(z)=\inf_{\alpha:E\to z,\ \alpha\subset Q}\int_\alpha g\,ds,
 \qquad u(z)=\min\{1,d(z)\},
\end{equation}
where the auxiliary paths are globally rectifiable. The infimum is $\infty$
if there is no finite-cost path. Constant paths at points of $E$ are allowed.

\begin{lemma}\label{lem:potential}
The functions $d$ and $u$ are lower semicontinuous. For every rectifiable
path $\beta\subset Q$ from $x$ to $y$,
\begin{equation}\label{eq:uppergradient}
 |u(x)-u(y)|\leq\int_\beta g\,ds.
\end{equation}
The actual pointwise function $u$ is absolutely continuous on almost every
horizontal and vertical line in $Q$, and $u\in W^{1,2}(\operatorname{int}Q)$.
\end{lemma}
\begin{proof}
We first justify the weighted-length compactness needed for measurability.
For a nonnegative continuous weight $a$ on $Q$ and a rectifiable curve
$\alpha:[0,1]\to Q$, consider the partition sums
\[
 S_P(\alpha)=\sum_i
 \min_{t\in[t_{i-1},t_i]}a(\alpha(t))\,
             |\alpha(t_i)-\alpha(t_{i-1})|.
\]
Each $S_P$ is continuous under uniform convergence of curves. Its supremum
over finite partitions is $\int_\alpha a\,ds$. Indeed, fine partitions make
the oscillation of $a\circ\alpha$ arbitrarily small, while further refinement
approximates the total variation of $\alpha$ by chord sums. The resulting
error is bounded by the weight oscillation times the curve length plus the
chord-variation error times $\max_Qa$. Thus weighted length for $a$ is lower
semicontinuous under uniform convergence.

A nonnegative lower semicontinuous $g$ on compact $Q$ is the increasing
pointwise limit of continuous nonnegative functions; explicitly one may use
\[
 a_m(x)=\min\left\{m,\inf_{y\in Q}\bigl(g(y)+m|x-y|\bigr)\right\}.
\]
These functions are continuous and increasing. If their limit at $x$ were
strictly below $g(x)$, bounded-cost near-minimizers would tend to $x$, contrary
to lower semicontinuity. This also proves convergence where $g(x)=\infty$.
Monotone convergence of line integrals now shows that weighted length for
$g$ is lower semicontinuous as well.

Suppose $z_j\to z$ and $L=\liminf_jd(z_j)<\infty$. Pass to a subsequence
with $d(z_j)\to L$, and choose paths $\alpha_j$ in~\eqref{eq:distance} of
cost at most $d(z_j)+1/j$. Since $g\geq1$, their Euclidean lengths are
bounded. Constant-speed parametrization on $[0,1]$ gives a common Lipschitz
bound. Arzel\`a--Ascoli supplies a uniformly convergent subsequence in compact
$Q$. Its limit $\alpha$ is rectifiable, starts in compact $E$, ends at $z$,
and satisfies
\[
 d(z)\leq\int_\alpha g\,ds
       \leq\liminf_j\int_{\alpha_j}g\,ds\leq L.
\]
Hence $d$ and its truncation $u$ are lower semicontinuous and Borel. This
argument also covers infinite distances without assuming that the infimum
is attained.

Concatenation with $\beta$ gives the distance triangle inequality whenever
the costs on its right side are finite. Reversing $x,y$ and truncating yields
\eqref{eq:uppergradient}. If one distance is infinite and the other finite,
any finite-cost connection would make both finite; if both are infinite,
both truncated values are $1$. Thus the same inequality holds in all cases.

By Fubini and Cauchy--Schwarz, $g$ is integrable on almost every coordinate
line in $Q$. On each such line~\eqref{eq:uppergradient} holds for every
subsegment. Applied to disjoint subintervals, it proves absolute continuity
of the actual function $u$, with the coordinate derivative bounded by $g$
almost everywhere. Slice integration by parts against a test function,
followed by Fubini, identifies these derivatives as weak derivatives.
Since $0\leq u\leq1$ and $|\nabla u|^2\leq2g^2$ almost everywhere,
$u\in W^{1,2}(\operatorname{int}Q)$.
\end{proof}

Now suppose that a finite-energy Borel density $\rho$ is admissible for
$\Gamma(E,F)$. Apply Lemma~\ref{lem:majorant} to its restriction to $Q$.
Every auxiliary path ending at $F\cap Q$ is a member of the original
connecting family. Its $g$-cost is therefore at least $1$. Constant paths
at $E$ have cost zero, so the potential has the exact point values
\begin{equation}\label{eq:plates}
 u=0\text{ on }E,\qquad u=1\text{ on }F\cap Q.
\end{equation}
These values follow from admissibility for every curve. No arbitrary
modification of a Sobolev representative on the real axis is being made.

Choose $\eta\in C_c^\infty(\operatorname{int}Q)$ equal to $1$ near
$[0,1]\times\{0\}$, and extend $w=\eta(1-u)$ by zero to $\R^2$.
Then $w\in W^{1,2}(\R^2)$ has compact support. Its actual restrictions to
almost every vertical line are absolutely continuous. By~\eqref{eq:plates}
their values at height zero give
\begin{equation}\label{eq:trace}
 f(x)=w(x,0)=\ind_E(x)\quad\text{for almost every }x\in\R.
\end{equation}
The exceptional vertical lines have one-dimensional measure zero. No
capacity statement about their intersection with $E$ is inferred.

\section{The critical trace forces length zero}

We give the required special cases of the standard flat trace and critical
characteristic-function rigidity principles; compare
\cite[Propositions 3.4 and 3.8]{DNPV} and
\cite[Corollaries 1--2 and Remark 3]{Brezis}.

\begin{lemma}\label{lem:trace}
For compactly supported $w\in W^{1,2}(\R^2)$, let $f$ be its almost-everywhere
vertical absolute-continuity trace. Then
\begin{equation}\label{eq:gagliardo}
 \int_\R\frac{\|f(\cdot+h)-f\|_{L^2(\R)}^2}{h^2}\,dh<\infty.
\end{equation}
\end{lemma}
\begin{proof}
First take $w$ smooth and compactly supported, with unitary horizontal
Fourier transform $W(\xi,y)=\widehat{w(\cdot,y)}(\xi)$, using
$e^{-ix\xi}$ as frequency convention. Integration of the derivative of
$|W|^2$ on $y>0$, followed by $2ab\leq a^2+b^2$, gives
\[
 |\xi|\,|W(\xi,0)|^2
 \leq2|\xi|\int_0^\infty|W|\,|\partial_yW|\,dy
 \leq\int_0^\infty\bigl(\xi^2|W|^2+|\partial_yW|^2\bigr)\,dy.
\]
Consequently Plancherel implies
\begin{equation}\label{eq:homtrace}
 \int_\R|\xi|\,|\widehat f(\xi)|^2\,d\xi
 \leq\int_{y>0}|\nabla w|^2\,dA.
\end{equation}
There is no division by $\xi$, including at zero frequency.

For a general compactly supported $w$, mollification gives compactly
supported smooth $w_j\to w$ in $W^{1,2}$. The traces converge to the actual
vertical trace in $L^2$. To see this directly, absolute continuity on almost
every vertical line gives, for each such compactly supported $v$,
\[
 |v(x,0)|^2\leq
 \int_0^\infty\bigl(|v(x,y)|^2+|\partial_yv(x,y)|^2\bigr)\,dy.
\]
This follows by integrating $\partial_y|v|^2$ and using $2|v||v'|\leq
|v|^2+|v'|^2$. Apply it to $v=w_j-w$ and integrate in $x$.
Plancherel gives $L^2$ convergence of the trace Fourier transforms; a
subsequence converges almost everywhere. Fatou then extends
\eqref{eq:homtrace} to $w$.

Finally Plancherel, Tonelli, and the substitution $t=\xi h$ give
\[
 \int_\R\frac{\|f(\cdot+h)-f\|_2^2}{h^2}\,dh
 =C\int_\R|\xi|\,|\widehat f(\xi)|^2\,d\xi,
 \qquad
 C=\int_\R\frac{|e^{it}-1|^2}{t^2}\,dt\in(0,\infty).
\]
The defining integrand for $C$ is bounded near zero and at most $4/t^2$
away from zero. This proves~\eqref{eq:gagliardo}.
\end{proof}

\begin{lemma}\label{lem:rigidity}
If a measurable $E\subset[0,1]$ has $\ind_E$ satisfying
\eqref{eq:gagliardo}, then $\length(E)=0$.
\end{lemma}
\begin{proof}
Write $f=\ind_E$. For $0<h<1$, put $N=\lceil2/h\rceil\leq3/h$.
Since $Nh\geq2$, the supports of $f$ and $f(\cdot+Nh)$ are disjoint.
Telescoping $N$ translations in $L^1$ yields
\[
 2\length(E)=\|f(\cdot+Nh)-f\|_1
 \leq N\|f(\cdot+h)-f\|_1.
\]
Characteristic-function differences satisfy $|f(x+h)-f(x)|^2=
|f(x+h)-f(x)|$. Hence
\[
 \|f(\cdot+h)-f\|_2^2\geq\tfrac23\length(E)h.
\]
The integral in~\eqref{eq:gagliardo}, restricted to $0<h<1$, is at least
$\frac23\length(E)\int_0^1dh/h$. Its finiteness forces $\length(E)=0$.
\end{proof}

Applying these lemmas to~\eqref{eq:trace} proves
\begin{equation}\label{eq:null}
 \rho\in\Adm\Gamma(E,F),\quad \int_{\R^2}\rho^2\,dA<\infty
 \quad\Longrightarrow\quad\length(E)=0.
\end{equation}
This argument includes positive-length sets without any interval. It gives
no capacity conclusion for the remaining null sets.

\section{Small circles at every endpoint}\label{sec:circles}

\begin{lemma}\label{lem:circles}
Let $E\subset[0,1]$ have length zero. If a Borel $\rho\in L^2(\R^2)$ is
admissible for $\Gamma(E,\R\setminus E)$, set
$q=\rho+\ind_{B(0,3)}$. Then
\begin{equation}\label{eq:infinitecost}
 \int_\gamma q\,ds=\infty
 \qquad\text{for every }\gamma\in\Gamma(E,S_2).
\end{equation}
\end{lemma}
\begin{proof}
The function $q$ is Borel, belongs to $L^2(\R^2)$, and remains admissible
for $\Gamma(E,F)$. Suppose a connector from a particular $e\in E$ to
$S_2$ has finite $q$-cost. Stop it at its first hit of $S_2$. Its preceding
image lies in $B(0,2)$, where $q\geq1$. Thus the lengths of all compact
interior subcurves are bounded by its finite $q$-cost. Continuity and the
endpoint limits show that the stopped curve has finite total variation,
including at the initial endpoint. It is therefore a globally rectifiable
curve and admits an arclength parametrization
$\gamma:[0,L]\to\overline{B(0,2)}$ after constant portions are removed.
Its line integral agrees with the increasing interior-subcurve integral:
endpoints have zero arclength measure. In particular $q\circ\gamma$ is
integrable on $[0,L]$. This explains why no infinite-length endpoint
behavior is lost in passing to a finite-cost curve. Its terminal point is
at distance at least $1$ from $e$.

Keep this actual endpoint $e$ fixed. Define the extended circle energy
\[
 A_e(r)=\int_{|z-e|=r}q(z)^2\,ds.
\]
Polar integration gives $\int_0^\delta A_e(r)\,dr<\infty$ for every finite
$\delta$. Also $e+r\notin E$ for almost every $r>0$, because
$\length(E)=0$. There are therefore radii $r_j\downarrow0$, $r_j<1$,
with
\begin{equation}\label{eq:radii}
 e+r_j\in F,\qquad r_jA_e(r_j)\longrightarrow0.
\end{equation}
For completeness, the radii with $e+r\in F$ and $A_e(r)<\infty$ have full
measure. If, for some $\varepsilon,\delta>0$, all such $0<r<\delta$
satisfied $rA_e(r)\geq\varepsilon$, then $\int_0^\delta A_e(r)\,dr$
would dominate $\varepsilon\int_0^\delta dr/r=\infty$. Thus arbitrarily
small good radii have arbitrarily small $rA_e(r)$, and a decreasing sequence
can be chosen recursively. This selection is made separately for each
fixed endpoint; no exceptional set of points of $E$ is discarded.

Let $t_j$ be the first arclength parameter with
$|\gamma(t_j)-e|=r_j$, and put $p_j=\gamma(t_j)$. Such parameters exist
since the terminal distance is at least $1$. Moreover $t_j\to0$.
Indeed, an arclength-parametrized curve is nonconstant on every initial
interval $[0,t]$ of positive length. Some point of that interval has
positive distance from $e$, so every sufficiently small disk centered at
$e$ is exited before $t$. Absolute continuity of the finite cost integral
therefore gives
\begin{equation}\label{eq:initialcost}
 \int_{\gamma|_{[0,t_j]}}q\,ds
 =\int_0^{t_j}q(\gamma(s))\,ds\longrightarrow0.
\end{equation}
Returns to $e$, retracing, or winding do not affect this first-exit argument.

Join $p_j$ to the real point $e+r_j\in F$ along a circular arc centered at
$e$. Cauchy--Schwarz on the full circle bounds the added cost by
\begin{equation}\label{eq:circlecost}
 \int_{|z-e|=r_j}q\,ds
 \leq\bigl(2\pi r_j A_e(r_j)\bigr)^{1/2}\longrightarrow0.
\end{equation}
The initial portion and this arc concatenate to an actual globally
rectifiable connector from $E$ to the original real complement $F$.
Interior plate contacts are allowed by definition. By
\eqref{eq:initialcost}--\eqref{eq:circlecost}, its $q$-cost is eventually
less than $1$, contradicting admissibility. This proves
\eqref{eq:infinitecost} for every connector, under either stated curve
convention.
\end{proof}

\begin{proof}[Proof of Theorem~\ref{thm:main}]
For empty $E$, the family to $S_2$ is empty and has modulus zero. Otherwise
suppose $\Mod\Gamma(E,F)<\infty$. By the definition of the infimum there
exists an admissible Borel $\rho$ of finite squared integral; no extremal
density is needed. Equation~\eqref{eq:null} gives $\length(E)=0$.
Lemma~\ref{lem:circles} supplies $q\in L^2(\R^2)$ whose cost on every
$E$-to-$S_2$ connector is infinite. Consequently $q/n$ is admissible for
the entire family for every positive integer $n$, and
\[
 0\leq\Mod\Gamma(E,S_2)
 \leq n^{-2}\int_{\R^2}q^2\,dA\longrightarrow0.
\]
Infinite-cost curves are included in this assertion, rather than removed
as an exceptional family. The contrapositive proves the theorem.
\end{proof}

\section{Context and status}

The critical trace and rigidity principles are classical
\cite{DNPV,Brezis}. Pointwise majorants, weighted-distance potentials,
and all-curve upper-gradient inequalities also occur in the established
framework of Heinonen--Koskela~\cite{HK}; no novelty of these mechanisms
is asserted. The present proof supplies the needed pointwise construction
and the separate null-set argument without importing a touching-plate
capacity identity or quasicontinuous boundary values.

The accompanying research audit also verifies a complete indirect
deduction assembled in October 2026 from older compact-condenser and
Sobolev theory, starting with compact-plate results such as
Hesse's~\cite{Hesse}. That assembly is not evidence that an earlier author
printed the exact resolution. No earlier explicit resolution was
authenticated in the bounded search, and historical priority remains
undetermined. In particular, the historical report in the 2018 source
does not establish present openness. The complete problem page in the
2019 Springer edition was unavailable for inspection, and the relevance
of a 1982 Goldstein--Jackson conference item entitled ``A problem of
A.~A.~Gon\v car'' remains unresolved. No firstness, current-openness, or
new-mechanism claim is made here.

\paragraph{Disclosure.}
AI systems were used extensively in solving, proof drafting, source
checking, and adversarial verification. This is an unrefereed research
note. Its internal mathematical audits are not external human peer review,
journal acceptance, or a formal proof-assistant certificate. The proof
uses basic measure theory, Arzel\`a--Ascoli, elementary Sobolev
mollification, and Plancherel; no computation is offered as a substitute
for the analytic argument.

\begin{thebibliography}{9}
\bibitem{Brezis}
H.~Brezis, \emph{How to recognize constant functions. Connections with
Sobolev spaces} (2002), Corollaries~1--2 and Remark~3,
\href{https://doi.org/10.1070/RM2002v057n04ABEH000533}{doi:10.1070/RM2002v057n04ABEH000533}.
\bibitem{DNPV}
E.~Di Nezza, G.~Palatucci, and E.~Valdinoci,
\emph{Hitchhiker's guide to the fractional Sobolev spaces},
arXiv:1104.4345v3 (19 November 2011), Propositions~3.4 and~3.8,
\url{https://arxiv.org/abs/1104.4345v3}.
\bibitem{HL}
W.~K.~Hayman and E.~F.~Lingham,
\emph{Research Problems in Function Theory (New Edition)},
arXiv:1809.07200v2 (21 September 2018), Problem~7.58,
printed p.~178, PDF p.~179,
\url{https://arxiv.org/abs/1809.07200v2}.
\bibitem{HK}
J.~Heinonen and P.~Koskela,
\emph{Quasiconformal maps in metric spaces with controlled geometry},
Acta Math. \textbf{180} (1998), 1--61,
\href{https://doi.org/10.1007/BF02392747}{doi:10.1007/BF02392747}.
\bibitem{Hesse}
J.~Hesse, \emph{A $p$-extremal length and $p$-capacity equality},
Ark. Mat. \textbf{13} (1975), 131--144, Theorem~5.5,
\href{https://doi.org/10.1007/BF02386202}{doi:10.1007/BF02386202}.
\end{thebibliography}
\end{document}
